\backhead{FURTHER EXERCISES}{khred}
\begin{exlist}
\item Using the fates of the sandy loam (0.638 transpired), find the irrigation needed to add 100 mm of transpiration.\exans{157 mm}
\item For a plateau of 0.3 and a dose scale of 4 mm, find the dose at which the drained share is half the plateau.\exans{2.8 mm}
\item A loop of gain 0.4 amplifies a first change of 5 per cent. Find the final change. What gain would double the first change?\exans{8.3 per cent; 0.5}
\item Irrigation raises the yield from 4.2 to 6.6 Mg\,ha$^{-1}$ with 180 mm. Find the irrigation water use efficiency.\exans{1.33 kg\,m$^{-3}$}
\item Nitrate leaching rises from 25 to 41 kg N\,ha$^{-1}$ when 150 mm of irrigation are added. Find the leaching per 100 mm of irrigation.\exans{10.7 kg N\,ha$^{-1}$}
\end{exlist}
\backhead{ESSAY QUESTIONS}{khred}
{\small\itshape Answer each question in about one page, with equations, a sketch and at least one number from this chapter.}\par\vspace{3pt}
\begin{exlist}
\item Discuss the three novel contributions of the study and the evidence for each.
\item Explain the fate of an irrigation millimetre and its use for advice.
\item Discuss the dose law and the rule of dosing against the daily use of the crop.
\item Explain the feedback loops of the study and how they could change under a warmer climate.
\item Discuss the coupling of irrigation and nitrogen leaching.
\item Write a short synthesis of the irrigation results for a farmer.
\end{exlist}
\dykband{EXERCISES IN SWAFL}
{\small Open SWAFL (desktop or online) with the benchmark project. The chapter of the \href{https://github.com/halilunaziru73-creator/Coupled-Water-Fluxes-Pore-Scale-Dynamics-Precision-Irrigation-Mulching-Feedback-Loops/raw/main/Software/v1.1.0/SWAFL_Desktop/docs/manual_v1.1.0/SWAFL_Manual_v1.1.0.pdf}{SWAFL manual} that explains each task is given in brackets; record your values in a table and compare them with what is expected (\emph{Check}).}\par\vspace{3pt}
\begin{exlist}
\item {\color{khblue}[Manual, Chapter 6]} Reproduce the irrigation of E2 for the loam at four triggers.\par\noindent{\small\itshape Check: irrigation falls with the trigger.}
\item {\color{khblue}[Manual, Chapter 12]} Reproduce the effect of dose on drainage in the drip layer.\par\noindent{\small\itshape Check: small doses drain less.}
\item {\color{khblue}[Manual, Chapter 10]} Reproduce the feedback loops of a warmer climate.\par\noindent{\small\itshape Check: the loops change sign in dry years.}
\item {\color{khblue}[Manual, Chapter 22]} Read the nitrogen budget of the irrigated and the rainfed field.\par\noindent{\small\itshape Check: irrigation raises leaching.}
\end{exlist}
